% !TeX root = ../analysis1.tex
\chapter{Infinite Series}\label{chap:analysis-series}

\section{Partial sums and convergence}
\begin{definition}[Infinite series]
For a sequence $(a_n)$, define its partial sums by
\[
S_N=\sum_{n=1}^N a_n.
\]
The series $\sum_{n=1}^{\infty}a_n$ converges to $S\in\mathbb R$ if
$S_N\to S$. Otherwise the series diverges. A series is therefore defined
through a limit of finite sums.
\end{definition}
\begin{example}[A warning about grouping]
The series $1-1+1-1+\cdots$ has partial sums $1,0,1,0,\ldots$, so it diverges.
Grouping it as $(1-1)+(1-1)+\cdots$ or as $1+(-1+1)+\cdots$ does not
assign a sum to the original divergent series.
\end{example}

\begin{proposition}[Basic properties]
Suppose $\sum a_n=S$ and $\sum b_n=T$. Then:
\begin{enumerate}
\item $a_n\to0$.
\item The tails satisfy $\sum_{n=k}^{\infty}a_n\to0$ as $k\to\infty$.
\item For $c\in\mathbb R$, $\sum ca_n=cS$.
\item $\sum(a_n+b_n)=S+T$.
\end{enumerate}
\end{proposition}
\begin{proof}
Since $a_n=S_n-S_{n-1}$ for $n\geq2$, the first claim follows from the
algebra of limits. The tail from $k$ is $S-S_{k-1}$, which tends to zero.
The last two claims follow by applying the constant-multiple and sum rules
to the corresponding finite partial sums. The tail argument and sum-rule
details complete exercises left to the reader in the source.
\end{proof}
\begin{remark}
The condition $a_n\to0$ is necessary, but not sufficient, for convergence.
The harmonic series below is the standard counterexample.
Changing, adding or deleting finitely many terms does not affect convergence,
although it can change the value of the sum.
\end{remark}

\section{Geometric and telescoping series}
\begin{example}[Geometric series]
For $r\ne1$,
\[
\sum_{n=0}^N r^n=\frac{1-r^{N+1}}{1-r}.
\]
If $|r|<1$, then $r^N\to0$, hence
\[
\sum_{n=0}^{\infty}r^n=\frac1{1-r},\qquad
\sum_{n=1}^{\infty}r^n=\frac r{1-r}.
\]
For completeness, if $0<|r|<1$, put $|r|^{-1}=1+h$, $h>0$.
Bernoulli's inequality $(1+h)^n\geq1+nh$ gives
$|r|^n\leq1/(1+nh)\to0$. Bernoulli's inequality follows by induction:
multiplying $1+nh$ by $1+h$ gives at least $1+(n+1)h$.
If $|r|\geq1$, the terms $r^n$ do not tend to zero, so the geometric
series diverges.
\end{example}

\begin{example}[Telescoping]
Using $1/[n(n+1)]=1/n-1/(n+1)$,
\begin{align*}
\sum_{n=1}^N\frac1{n(n+1)}
&=\sum_{n=1}^N\left(\frac1n-\frac1{n+1}\right)\\
&=1-\frac1{N+1}\longrightarrow1.
\end{align*}
Thus $\sum_{n=1}^{\infty}1/[n(n+1)]=1$.
\end{example}

\section{Comparison and absolute convergence}
For nonnegative terms $a_n\geq0$, the partial sums are increasing.
They either are bounded and converge, or are unbounded and tend to $+\infty$.
This observation is the basis of comparison tests.

\begin{theorem}[Comparison test]
If $0\leq a_n\leq b_n$ for all sufficiently large $n$ and $\sum b_n$
converges, then $\sum a_n$ converges. If instead $\sum a_n$ diverges, then
$\sum b_n$ diverges. When the inequalities hold for every $n$,
$\sum a_n\leq\sum b_n$ in the convergent case.
\end{theorem}
\begin{proof}
First suppose the inequalities hold for every $n$. The partial sums
$A_N=\sum_{n=1}^N a_n$ are increasing and satisfy
$A_N\leq B_N\leq\sum b_n$. Thus they converge by monotone convergence.
The sum inequality follows by passing to limits.
If comparison holds only after a finite index, apply this argument to the
tails, since finitely many initial terms do not affect convergence.
The divergence claim is the contrapositive.
\end{proof}
This includes the source's ``improved comparison test'': eventual comparison
is sufficient.

\begin{example}
Since $0\leq2^n/(3^n+1)\leq(2/3)^n$, the series
$\sum 2^n/(3^n+1)$ converges. Also, for $n\geq2$,
\[
0<\frac1{n^3-2}\leq\frac2{n^3}.
\]
Once convergence of $\sum1/n^3$ has been established, comparison proves
convergence of $\sum_{n=1}^{\infty}1/(n^3-2)$. Its first term is $-1$;
comparison is applied only from $n=2$ onward.
\end{example}

\begin{definition}[Absolute and conditional convergence]
A series $\sum a_n$ is \emph{absolutely convergent} if $\sum|a_n|$ converges.
It is \emph{conditionally convergent} if $\sum a_n$ converges but
$\sum|a_n|$ diverges.
\end{definition}

\begin{theorem}
Every absolutely convergent real series converges.
\end{theorem}
\begin{proof}
Put $a_n^+=\max\{a_n,0\}$ and $a_n^-=\max\{-a_n,0\}$.
Then $a_n=a_n^+-a_n^-$, $|a_n|=a_n^++a_n^-$, and
$0\leq a_n^\pm\leq|a_n|$. Both $\sum a_n^+$ and $\sum a_n^-$ converge
by comparison. Subtracting their partial-sum limits proves convergence
of $\sum a_n$.
\end{proof}

\begin{theorem}[$p$-series]
For $p>0$, the series $\sum_{n=1}^{\infty}n^{-p}$ converges if and only
if $p>1$.
\end{theorem}
\begin{proof}
For $0<p\leq1$, $n^{-p}\geq1/n$. The harmonic partial sums obey
\[
\sum_{n=1}^{2^M}\frac1n
=1+\sum_{j=1}^M\sum_{n=2^{j-1}+1}^{2^j}\frac1n
\geq1+\frac M2.
\]
Each block contains $2^{j-1}$ terms, each at least $2^{-j}$.
Thus the harmonic series, and consequently each series with $0<p\leq1$,
diverges to $+\infty$.

For $p>1$, group terms with $2^j\leq n<2^{j+1}$:
\[
\sum_{n=2^j}^{2^{j+1}-1}\frac1{n^p}
\leq\frac{2^j}{(2^j)^p}=\bigl(2^{1-p}\bigr)^j.
\]
The geometric series with ratio $2^{1-p}<1$ gives a common bound for all
partial sums. Since the partial sums increase, they converge.
\end{proof}
For irrational $p$, real powers are defined later using the exponential
and logarithm. The same estimates apply to those powers.

\begin{theorem}[Limit comparison]
Let $a_n,b_n>0$ eventually and suppose $a_n/b_n\to\ell$, where
$0<\ell<\infty$. Then $\sum a_n$ converges if and only if $\sum b_n$
converges.
\end{theorem}
\begin{proof}
Eventually $\ell/2<a_n/b_n<3\ell/2$, so compare both ways with positive
constant multiples of $b_n$. This proof supplies the tutorial result
cited, without proof, in the source.
\end{proof}

\section{Alternating series}
\begin{theorem}[Alternating series test]
If $a_n\geq0$, $a_{n+1}\leq a_n$, and $a_n\to0$, then
$\sum_{n=1}^{\infty}(-1)^{n+1}a_n$ converges.
\end{theorem}
\begin{proof}
Let $S_N=\sum_{n=1}^N(-1)^{n+1}a_n$. The even partial sums increase:
\[
S_{2N+2}-S_{2N}=a_{2N+1}-a_{2N+2}\geq0.
\]
The odd partial sums decrease:
$S_{2N+1}-S_{2N-1}=-a_{2N}+a_{2N+1}\leq0$.
Also $S_{2N}\leq S_{2N-1}\leq a_1$ and $S_{2N}\geq0$.
Thus both subsequences converge, and their limits agree because
$S_{2N-1}-S_{2N}=a_{2N}\to0$. Every sufficiently late partial sum is
an even or odd term close to that common limit, so the entire sequence
of partial sums converges.
\end{proof}
\begin{example}
The alternating harmonic series
$1-1/2+1/3-1/4+\cdots$ converges by this test, but not absolutely because
$\sum1/n$ diverges. Its sum is $\log2$; that identity is mentioned in
the source but not established by the alternating series test alone.
\end{example}

\section{Ratio and root tests}
\begin{theorem}[Ratio test]
Suppose $a_n\ne0$ eventually and
$|a_{n+1}|/|a_n|\to\ell$. Then:
\begin{enumerate}
\item If $\ell<1$, $\sum a_n$ converges absolutely.
\item If $\ell>1$, $|a_n|\to+\infty$, so $\sum a_n$ diverges.
\item If $\ell=1$, the test gives no conclusion.
\end{enumerate}
The divergence conclusion also holds when the ratio tends to $+\infty$.
\end{theorem}
\begin{proof}
For $\ell<1$, choose $r$ with $\ell<r<1$. Eventually
$|a_{n+1}|\leq r|a_n|$, and induction gives
$|a_{N+k}|\leq|a_N|r^k$. Comparison with a geometric series proves
absolute convergence.
For $\ell>1$, choose $1<r<\ell$. Eventually the reverse bound gives
$|a_{N+k}|\geq|a_N|r^k\to+\infty$.
\end{proof}

\begin{example}[The exponential series]\label{ex:analysis-exp}
Define, for every real $x$,
\[
\exp(x)=\sum_{n=0}^{\infty}\frac{x^n}{n!},\qquad 0!=1.
\]
For $x\ne0$, the ratio of successive absolute terms is $|x|/(n+1)\to0$,
so the series converges absolutely. For $x=0$, the sum is just $1$.
This defines $\exp$ before its continuity or differentiability is known.
\end{example}

\begin{theorem}[Root test]
Suppose $\sqrt[n]{|a_n|}\to\ell$. If $\ell<1$, $\sum a_n$ converges
absolutely. If $\ell>1$, it diverges. If $\ell=1$, the test is inconclusive.
\end{theorem}
\begin{proof}
For $\ell<1$, choose $\ell<r<1$. Eventually $|a_n|\leq r^n$, so comparison
proves absolute convergence. For $\ell>1$, choose $1<r<\ell$.
Eventually $|a_n|\geq r^n$, so the terms do not tend to zero.
\end{proof}

\begin{example}[Recognising inconclusive cases]
For both $a_n=1/n$ and $a_n=1/n^2$, the successive ratios tend to $1$.
The first series diverges and the second converges, so ratio limit $1$
cannot decide convergence.
\end{example}
\begin{example}[Using the root test]
For $a_n=(-1)^n((2n+7)/(5n+2))^n$,
\[
\sqrt[n]{|a_n|}=\frac{2n+7}{5n+2}\longrightarrow\frac25<1.
\]
Hence $\sum a_n$ converges absolutely.
\end{example}

\subsection{Choosing a convergence test}
First check whether $a_n\to0$. If not, the series diverges.
For nonnegative terms, try comparison or limit comparison with a geometric
or $p$-series. Factorials often suggest the ratio test; an expression raised
to the $n$th power often suggests the root test. For terms with changing
signs, test absolute convergence, or use the alternating series test when
its monotonicity and zero-limit hypotheses hold.

\section{Rearrangements}
\begin{definition}
A rearrangement of $\sum_{n=1}^{\infty}a_n$ is
$\sum_{n=1}^{\infty}a_{\pi(n)}$, where
$\pi:\mathbb N\to\mathbb N$ is a bijection. Every original term appears
exactly once.
\end{definition}

\begin{theorem}[Absolute convergence and rearrangement]
Every rearrangement of an absolutely convergent series converges to the
same sum.
\end{theorem}
\begin{proof}[Supplementary proof]
Let $S=\sum a_n$ and choose $K$ with $\sum_{n>K}|a_n|<\varepsilon$.
For sufficiently large $N$, the indices $\pi(1),\ldots,\pi(N)$ include
$1,\ldots,K$. After removing the terms common to the rearranged partial
sum and $S$, only some terms with indices greater than $K$ remain. Thus
\[
\left|S-\sum_{j=1}^N a_{\pi(j)}\right|
\leq\sum_{n>K}|a_n|<\varepsilon.
\]
The source states this theorem without its proof.
\end{proof}

\begin{theorem}[Riemann rearrangement theorem]
If a real series converges conditionally, then for every $L\in\mathbb R$
it has a rearrangement converging to $L$. There are also rearrangements
whose partial sums tend to $+\infty$ or $-\infty$.
\end{theorem}
\begin{proof}[Supplementary proof outline]
The positive terms have sum $+\infty$ and the negative terms have sum
$-\infty$; otherwise convergence of their difference would force absolute
convergence. To obtain a finite target $L$, take unused positive terms
in their original order until the partial sum exceeds $L$, then unused
negative terms until it falls below $L$, and repeat. Each crossing error
is bounded by the last term added, which tends to zero. Both lists are
exhausted, and any zero terms can be inserted, so this is a rearrangement
converging to $L$.

For $+\infty$, at stage $k$ insert the next unused negative term, then
enough unused positive terms to make the sum exceed $k$. At least one
positive term is used at each stage. The negative terms tend to zero,
so the intervening drops also tend to zero, and the partial sums tend to
$+\infty$. Reverse the roles to obtain $-\infty$.
\end{proof}

\begin{example}[Rearranging the alternating harmonic series]
Let $S=1-1/2+1/3-1/4+\cdots$. Taking two positive terms for each negative
term gives
\[
1+\frac13-\frac12+\frac15+\frac17-\frac14+\cdots.
\]
Its sum is $3S/2$, not $S$. To justify this, at the end of its $N$th block
its sum is
\[
T_{3N}=\sum_{k=1}^{2N}\frac1{2k-1}-\sum_{k=1}^N\frac1{2k}
=S_{4N}+\frac12\sum_{k=N+1}^{2N}\frac1k.
\]
The last sum equals $S_{2N}$, since both equal $H_{2N}-H_N$ with
$H_n=\sum_{k=1}^n1/k$. Hence $T_{3N}\to S+S/2$.
The one or two extra terms between block endpoints tend to zero, proving
the claimed limit for all partial sums.
\end{example}

\section{Cauchy products and the exponential addition law}
\begin{definition}[Cauchy product]
For series indexed from $0$, define
\[
c_n=\sum_{k=0}^n a_kb_{n-k}.
\]
The series $\sum_{n=0}^{\infty}c_n$ is their Cauchy product. Its terms
collect products along the diagonals $k+(n-k)=n$.
\end{definition}

\begin{theorem}
If $\sum_{n=0}^{\infty}a_n=A$ and $\sum_{n=0}^{\infty}b_n=B$ both
converge absolutely, their Cauchy product converges absolutely and has
sum $AB$.
\end{theorem}
\begin{proof}[Supplementary proof]
Let $A_* =\sum|a_n|$ and $B_* =\sum|b_n|$. For each $N$,
\[
\sum_{n=0}^N|c_n|
\leq\sum_{k+j\leq N}|a_k||b_j|\leq A_*B_*.
\]
Thus the product series converges absolutely.
Compare its $N$th partial sum with the rectangular product
\[
\left(\sum_{k=0}^N a_k\right)\left(\sum_{j=0}^N b_j\right).
\]
Their difference consists of terms
with $0\leq k,j\leq N$ and $k+j>N$. Each such term has either $k>N/2$
or $j>N/2$. Its total absolute value is therefore at most
\[
B_*\sum_{k>N/2}|a_k|+A_*\sum_{j>N/2}|b_j|\longrightarrow0.
\]
The rectangular product tends to $AB$, proving the result. The source
omits this proof.
\end{proof}

\begin{corollary}[Exponential addition law]\label{thm:analysis-exp-addition}
For all $x,y\in\mathbb R$, $\exp(x+y)=\exp(x)\exp(y)$.
\end{corollary}
\begin{proof}
Both exponential series converge absolutely. Their $n$th Cauchy-product
coefficient is, by the finite binomial theorem,
\[
\sum_{k=0}^n\frac{x^ky^{n-k}}{k!(n-k)!}
=\frac1{n!}\sum_{k=0}^n\binom nk x^ky^{n-k}
=\frac{(x+y)^n}{n!}.
\]
Summing gives the identity.
\end{proof}
In particular, $\exp(0)=1$ and $\exp(x)\exp(-x)=1$.
Moreover $\exp(x)>0$: it is positive for $x\geq0$ by its series, and for
$x<0$ it is the reciprocal of the positive number $\exp(-x)$.
